% Section 8 — Chronoceptive Calibration Error (compressed)

\section{The Chronoceptive Calibration Error $\cce$}
\label{sec:epsilon}

We aggregate the three axes into a single scalar for cross-agent comparison and the deployment bound of \S\ref{sec:agentic-timeline}.

\begin{definition}[Chronoceptive Calibration Error]
\label{def:cce}
$\cce(A) := \frac{1}{|\mathcal{A}|}\sum_{k \in \mathcal{A}} \mathrm{score}_k$, where $\mathcal{A}$ is the set of axes with at least one defined trajectory and each axis contributes the median of its per-trajectory error: $\mathrm{score}_{T_2} = \mathrm{med}\,|\CAR - 1|$ on T2.3, $\mathrm{score}_{T_3} = \mathrm{med}\,|\confab|$ on T3.1. $\cce = 0$ is grounded; $\ccestar = 0.20$ is the Augustine threshold. An axis with no trajectories is dropped from $\mathcal{A}$, never imputed as zero, so an agent is never credited for a measurement that was not taken.
\end{definition}

\paragraph{Which axes, and why medians.} Under this protocol $\mathcal{A} = \{T_2, T_3\}$. The $\twall$-axis coefficient $\alpha$ requires a wall-clock budget on a $\twall$-axis capability, and no capability here carries one: T2.3 is the only budgeted capability and it loads on $\tstep$. We report the two axes we measured rather than inflating the aggregate with a third we did not. Each axis is a median rather than a mean because two trajectories in the corpus are harness artefacts, not agent behaviour --- gpt-4o-mini on T2.3 with $\twall = 14.9$\,h against a $900$\,s budget, a stalled request rather than fifteen hours of inference. Under a mean, those two alone move that agent's $\cce$ from $1.06$ to $1.99$ and hand it last place.

\paragraph{A correction to the convention.} An earlier draft of this table scored an unmeasured axis as $0$ while leaving its weight in the denominator, which credited every agent with a free perfect third --- and credited o4-mini, whose original corpus had no T2.3 cell, with two. Every $\cce$ below follows the definition above. The correction raises all $\cce$ uniformly for agents with a full corpus, so it strengthens rather than rescues the claim that no agent crosses $\ccestar$; it changes the ordering in exactly one place, moving o4-mini from mid-panel to the bottom, where its $|\confab| = 1.54$ belongs. The o4-mini reasoning-effort ladder is reported on the $\confab$ axis alone at all three tiers, since T2.3 was only run at the default effort.

\begin{figure}[h]
\centering
\includegraphics[width=\linewidth]{figures/epsilon_panel.pdf}
\caption{$\cce$ over the 13 scored agents. No agent crosses $\ccestar$; the closest, Sonnet 4.6, sits $2.7\times$ above. $\mathrm{score}_{T_2}$ lies in $[0.92, 1.00]$ on every agent --- irreducible by any narrative-axis intervention. DeepSeek-R1-Distill-14B is the 14th panel configuration and is unscored (it declines to spend its reasoning budget). Full sub-capability decomposition in App.~\ref{app:l3-detail}.}
\label{fig:epsilon-panel}
\end{figure}

\paragraph{Two findings.} \textbf{(1) The action axis is binding.} $\mathrm{score}_{T_2}$ lies in $[0.92, 1.00]$ on every agent in the panel, Setting~A and Setting~B alike. Projecting $\mathrm{score}_{T_3} \to 0$ --- which would require crossing both Theorem~\ref{thm:cit} and Theorem~\ref{thm:reverse-scaling} --- still leaves $\cce \geq 0.46$, a factor of $2.3$ above $\ccestar$. No narrative-axis intervention can reach the threshold, because the narrative axis is not what is holding the agent back. \textbf{(2) Reverse-Scaling operates on the aggregate.} Adding extended thinking to Sonnet 4.6 moves $\cce$ the wrong way, $0.536 \to 0.566$ under Setting~A and $0.537 \to 0.585$ under Setting~B, so the thinking variant is no longer the best agent in the panel. The o4-mini effort ladder moves it further: on the $\confab$ axis, $1.23 \to 1.54 \to 1.67$ from low to high effort. More reasoning compute moves the agent away from the threshold on the aggregate, exactly as it does on the axis Theorem~\ref{thm:reverse-scaling} constrains.

We flag the size of the first effect against the second. The $\pm$ thinking contrast is a $5.6\%$ move on Setting~A and rests on a thinking-variant $\confab$ cell with $n_\rho = 10$; it is consistent with Reverse-Scaling on the aggregate but is not on its own strong evidence for it. The effort ladder, at $n_\rho = 16/41/17$ across three tiers and a $36\%$ total move, is.

This is the paper's central empirical claim. The next section turns it into a deployment bound.
